\documentclass[12pt,a4paper]{article}

\usepackage{cmap}
\usepackage[utf8]{inputenc}
\usepackage[T1]{fontenc}
\usepackage{lmodern}
\usepackage[brazil]{babel}
\usepackage{geometry}
\usepackage{graphicx}
\usepackage{booktabs}
\usepackage{array}
\usepackage{microtype}
\usepackage[hidelinks]{hyperref}
\usepackage{csquotes}
\usepackage[backend=biber,style=numeric,sorting=none]{biblatex}

\AtEveryBibitem{%
  \clearfield{note}%
  \clearfield{file}%
}
\addbibresource{../references.bib}

\geometry{margin=2.5cm}
\setlength{\emergencystretch}{2em}

\newcommand{\cooe}{CO$_2$e}
\newcommand{\ttw}{TTW}
\newcommand{\wtw}{WTW}

\title{CabotageLens: framework computacional auditável para comparação porta a porta entre rodovia e cabotagem no Brasil}
\author{Felipe de Sá Proença}
\date{\today}

\begin{document}

\maketitle

\begin{abstract}
Comparações entre transporte rodoviário e cabotagem exigem uma unidade funcional comum, a representação completa dos acessos terrestres e uma fronteira ambiental consistente. Este artigo apresenta o CabotageLens, um framework computacional auditável para comparar uma rota rodoviária direta com uma cadeia rodoviária--cabotagem--rodoviária em corredores brasileiros. O modelo marítimo combina registros de carga, atracação e sequência portuária da ANTAQ com intensidades de combustível do EU MRV associadas pelo número IMO. Para cada par portuário direcional, são utilizadas todas as viagens nas quais a origem ocorre antes do destino, incluindo ligações diretas, viagens com escalas intermediárias e diferentes corredores observados. A carga a bordo é reconstruída em cada subtrecho e define, com a distância navegada, o trabalho de transporte da viagem. A intensidade de cada navio é obtida prioritariamente pelo IMO; quando essa correspondência não está disponível, aplica-se uma estatística robusta de classe ou tipo de navio, com aparagem simétrica de 1\% das caudas e proveniência explícita. A intensidade representativa do par é a mediana das intensidades das viagens ponderada pelo trabalho de transporte, enquanto a escolha do corredor determina apenas a sequência portuária e a distância aplicadas ao cenário. A base contém 1.324 viagens observadas, 389 IMOs e 177 pares direcionais. No par Santos--Manaus, 89 observações viagem--origem--destino distribuídas em 22 corredores resultam em intensidade de 9,322050~g/(t$\cdot$nm). O framework preserva a separação entre evidência observada, fallback estatístico e parâmetros de cenário, e limita sua interpretação a emissões operacionais \ttw{} de \cooe{} e a um proxy de custo operacional modelado.

\textbf{Palavras-chave}: cabotagem; transporte rodoviário; transporte multimodal; ANTAQ; EU MRV; emissões operacionais; logística; Brasil.
\end{abstract}

\section{Introdução}
\label{sec:introducao}

O transporte de cargas no Brasil permanece fortemente dependente do modo rodoviário. Essa configuração oferece capilaridade territorial e flexibilidade operacional, mas concentra fluxos de longa distância em um modo sensível ao consumo de diesel, à infraestrutura terrestre e à volatilidade de custos. Para corredores conectáveis à costa, a cabotagem constitui uma alternativa técnica relevante, especialmente em um país com extensa faixa litorânea e grandes centros produtores e consumidores próximos a complexos portuários \cite{icct2022}.

A relevância da cabotagem não implica superioridade automática sobre o transporte rodoviário. A pergunta técnica não é se um navio é melhor que um caminhão em abstrato, mas como duas cadeias capazes de transportar a mesma remessa entre a mesma origem e o mesmo destino se comparam. O resultado depende dos portos utilizados, das distâncias terrestres e marítimas, da utilização da embarcação, da carga alocada ao transporte, das interfaces portuárias e da fronteira ambiental e econômica adotada \cite{shortsea2019,modalshiftreview2020}. Comparar uma perna marítima isolada com uma viagem rodoviária completa produziria uma assimetria metodológica.

O CabotageLens foi concebido para realizar essa comparação de forma porta a porta. A alternativa rodoviária direta e a alternativa rodoviária--cabotagem--rodoviária são avaliadas para a mesma carga, com decomposição por perna, fontes registradas e avisos de qualidade. O framework também separa emissões operacionais \ttw{}, emissões de ciclo de vida e resultados monetários. O custo calculado representa um proxy operacional modelado, e não frete comercial, tarifa contratada ou demonstração de viabilidade logística \cite{competitiveness2024,decarb2024,maritimelca2024}.

A contribuição específica deste artigo é apresentar uma formulação marítima baseada em atividade observada. Registros da ANTAQ permitem reconstruir viagens, paradas e cargas a bordo, enquanto o EU MRV fornece indicadores anuais de combustível por trabalho de transporte associados aos navios por IMO. Essa combinação torna possível representar ligações portuárias direcionais sem fixar antecipadamente um corredor e sem atribuir a toda a frota um único navio médio.

\section{Revisão da literatura e fundamentação metodológica}
\label{sec:contexto}

A literatura sobre cabotagem brasileira e sobre o programa BR do Mar sustenta a relevância do modo para corredores longos, mas não substitui a avaliação de relações origem--destino específicas \cite{icct2022}. Estudos de competitividade em rede mostram que frequência, tempo, estoque, terminais, confiabilidade, risco e disponibilidade de serviço participam da decisão modal real \cite{competitiveness2024}. O CabotageLens concentra-se na comparação técnica de rotas, combustível, emissões operacionais e proxy de custo, sem representar integralmente essas variáveis comerciais.

Na literatura de \textit{short sea shipping}, a direção e a magnitude do desempenho ambiental dependem da distância, dos acessos terrestres, do tipo de navio, da utilização, dos fatores de emissão e da forma de alocação à carga \cite{shortsea2019}. Revisões de mudança modal reforçam que uma vantagem energética potencial não implica adoção automática do modo marítimo \cite{modalshiftreview2020}. A unidade funcional deve, portanto, abranger a remessa completa, e não veículos isolados.

O EU MRV oferece dados reportados por navio e indicadores normalizados por distância e por trabalho de transporte. Esses indicadores permitem relacionar combustível ao serviço prestado e são mais adequados à alocação de uma carga do que o consumo anual absoluto da embarcação \cite{eumrv2025}. No CabotageLens, o indicador empregado é expresso em g/(t$\cdot$nm), preservando a associação por IMO e o ano de referência. A base da ANTAQ fornece a contrapartida operacional brasileira: identificação do navio, sequência de atracações e movimentos de carga em cabotagem conteinerizada \cite{antaq2025}.

A fronteira ambiental adotada é a de emissões operacionais \ttw{} de \cooe{}. Fontes \wtw{}, LCA ou exclusivamente de CO$_2$ são úteis para discutir limites e trabalhos futuros, mas não são tratadas como fatores diretamente intercambiáveis com a saída operacional do modelo \cite{decarb2024,maritimelca2024}. Operações portuárias e consumo de navios atracados também requerem atenção, pois dependem de terminal, produtividade, classe do navio e tempo de permanência \cite{berth2009,berthairquality2010,shipops2022}.

\begin{table}[htbp]
\centering
\small
\caption{Fronteiras metodológicas do CabotageLens.}
\label{tab:fronteiras}
\begin{tabular}{p{0.23\textwidth}p{0.34\textwidth}p{0.35\textwidth}}
\toprule
Dimensão & Incluído & Fora da fronteira \tabularnewline
\midrule
Alternativas & Rodovia direta e cadeia rodoviária--cabotagem--rodoviária porta a porta & Comparação entre navio isolado e caminhão para a viagem completa \tabularnewline
Emissões & Emissões operacionais \ttw{} de \cooe{} por remessa & \wtw{}, LCA, fabricação de ativos e inventário completo de poluentes locais \tabularnewline
Custo & Proxy de custo operacional modelado & Frete comercial, negociação, seguro, estoque, demurrage, detention e booking \tabularnewline
Dados marítimos & Viagens ANTAQ, intensidades EU MRV e fallbacks identificados & Interpretação de fallback como medição individual do navio \tabularnewline
Serviço & Corredores observados na janela de dados & Garantia de frequência, slot, escala futura ou disponibilidade comercial \tabularnewline
\bottomrule
\end{tabular}
\end{table}

\section{Metodologia}
\label{sec:metodologia}

\subsection{Unidade funcional e alternativas}
\label{subsec:unidade-funcional}

A unidade funcional é o transporte de uma quantidade especificada de carga conteinerizada entre uma origem e um destino no Brasil. A configuração de referência utiliza 1~TEU e 14~t por remessa. A massa permanece um dado de cenário e pode ser alterada pelo usuário; ela não deve ser confundida com a carga reconstruída a bordo dos navios da ANTAQ.

A alternativa rodoviária direta contém uma rota terrestre origem--destino. A alternativa multimodal contém o acesso rodoviário ao porto de origem, as interfaces portuárias representáveis, a perna marítima, as interfaces no porto de destino e o acesso rodoviário final. Todos os componentes são agregados para a mesma remessa \cite{shortsea2019,competitiveness2024}.

\subsection{Reconstrução das viagens e da carga a bordo}
\label{subsec:reconstrucao-antaq}

O processamento da ANTAQ combina as tabelas de Carga, Atracação e Tempos de Atracação. Os registros de carga informam movimentos embarcados e desembarcados, massa, quantidade e portos associados; os registros de atracação fornecem porto, tempo, embarcação e IMO. Chamadas consecutivas pertencentes ao mesmo complexo portuário canônico são colapsadas após a soma dos movimentos de carga, de modo que a normalização dos nomes não descarte atividade observada.

Para cada IMO, as escalas são ordenadas e segmentadas em viagens. Em cada parada, a variação líquida de carga atualiza o inventário a bordo. Quando a primeira escala observada implicaria saldo negativo, reconstrói-se a carga inicial mínima necessária para manter a série não negativa. A carga associada ao trecho que sai da parada $s$, denotada por $m_{v,s}$, é o saldo acumulado após os movimentos dessa parada. Esse procedimento produz uma sequência de subtrechos contíguos, cada qual com porto inicial, porto final, carga a bordo e fonte de distância.

Para um par ordenado de portos $(o,d)$, a viagem $v$ é elegível quando $o$ aparece antes de $d$ em sua sequência. O recorte inclui viagens diretas e viagens com qualquer número de escalas intermediárias. Todas as sequências completas observadas entre os dois portos são avaliadas, inclusive as que percorrem corredores diferentes. Um corredor nunca é sintetizado pela junção de subtrechos pertencentes a viagens distintas. Cada viagem contribui uma vez para o mesmo par ordenado; se chamadas repetidas gerarem mais de um recorte possível dentro da viagem, utiliza-se primeiro o recorte direto e, na sua ausência, o recorte completo de menor distância.

\subsection{Intensidade marítima por IMO e fallback robusto}
\label{subsec:intensidade-mrv}

A intensidade $I_v$ atribuída à viagem segue uma hierarquia de fontes. A primeira opção é o último indicador positivo de combustível por trabalho de transporte encontrado para o IMO no EU MRV. Essa correspondência preserva a identidade da embarcação e não sofre aparagem ou substituição, mesmo quando o valor se encontra nas caudas da distribuição.

Na ausência de correspondência positiva por IMO, utiliza-se uma estatística de fallback. Uma classe de navio explicitamente disponível tem prioridade sobre o tipo. Para a classe, emprega-se a média aparada em 1\% por cauda registrada no artefato de eficiência; a mediana da classe é usada somente quando essa estatística não está disponível. Para o tipo, os últimos indicadores positivos de cada IMO são ordenados, removem-se $\lfloor 0{,}01n\rfloor$ observações de cada cauda e calcula-se a média dos valores retidos. Se a amostra for pequena demais para excluir ao menos uma observação em cada cauda, utiliza-se a mediana. No recorte conteinerizado sem metadado de tipo, \textit{container ship} é o tipo documentado.

A regra de outliers incide, portanto, sobre a construção do fallback coletivo, não sobre observações exatas associadas ao IMO. Para cada intensidade, o sistema registra a fonte, o nível da associação, o estimador, a regra de aparagem, o tamanho bruto e retido da amostra, a quantidade excluída e os limites retidos. Essa proveniência diferencia explicitamente uma observação específica do navio de uma imputação estatística de classe ou tipo.

\subsection{Trabalho de transporte e intensidade representativa do par}
\label{subsec:intensidade-par}

Seja $\mathcal{S}_{v,o,d}$ o conjunto de subtrechos contíguos da viagem $v$ entre a origem portuária $o$ e o destino portuário $d$. Para cada subtrecho $s$, $d_{v,s}$ é a distância em milhas náuticas e $m_{v,s}$ é a carga reconstruída a bordo em toneladas. O trabalho de transporte observado da viagem é

\begin{equation}
W_{v,o,d}=\sum_{s\in\mathcal{S}_{v,o,d}}m_{v,s}\,d_{v,s}.
\label{eq:trabalho-transporte}
\end{equation}

O combustível atribuído à atividade conteinerizada observada entre esses portos é

\begin{equation}
F_{v,o,d}^{\mathrm{obs}}=\frac{I_v}{1000}
\sum_{s\in\mathcal{S}_{v,o,d}}m_{v,s}\,d_{v,s},
\label{eq:combustivel-observado}
\end{equation}

em que $I_v$ é expresso em g/(t$\cdot$nm) e $F_{v,o,d}^{\mathrm{obs}}$ em kg. A soma abrange todos os subtrechos existentes entre origem e destino; assim, uma viagem com escalas intermediárias não é reduzida a uma ligação direta hipotética.

A intensidade aplicada a um cenário não é escolhida com base em um único corredor. Todas as observações viagem--origem--destino resolvidas para o mesmo par participam diretamente do estimador. Ordenando as intensidades $I_v$ e usando $W_{v,o,d}$ como peso, define-se a mediana ponderada inferior:

\begin{equation}
I_{o,d}^{\mathrm{rep}}=
\min\left\{x:\sum_{v:I_v\leq x}W_{v,o,d}
\geq\frac{1}{2}\sum_vW_{v,o,d}\right\}.
\label{eq:mediana-ponderada}
\end{equation}

Esse estimador representa todas as viagens elegíveis do par, sem calcular primeiro uma média por corredor. Observações com trabalho de transporte nulo não influenciam a mediana quando existe peso positivo. Se todas as observações resolvidas tiverem trabalho nulo, utiliza-se uma mediana simples com fonte e método específicos. A mediana ponderada reduz a influência de caudas elevadas sem selecionar um valor em função de um resultado desejado.

\subsection{Separação entre intensidade e corredor}
\label{subsec:corredor}

Os corredores observados cumprem duas funções distintas. A população completa de viagens determina $I_{o,d}^{\mathrm{rep}}$; a regra de seleção de corredor determina somente a sequência portuária e a distância do cenário. Entre alternativas completas e utilizáveis, dá-se preferência a um corredor direto observado. Se não houver ligação direta utilizável, escolhe-se o corredor completo de menor distância. Subtrechos de viagens diferentes não são combinados.

Para uma carga de cenário $M$ e um corredor selecionado $\mathcal{S}^{*}_{o,d}$, o consumo marítimo de navegação é

\begin{equation}
F_{o,d}^{\mathrm{cen}}=\frac{I_{o,d}^{\mathrm{rep}}\,M}{1000}
\sum_{s\in\mathcal{S}^{*}_{o,d}}d_s.
\label{eq:combustivel-cenario}
\end{equation}

As intensidades observadas no corredor selecionado são mantidas como diagnóstico, mas não substituem o valor representativo do par. A distância de cada subtrecho utiliza a matriz marítima quando há valor positivo; na falta desse dado, o sistema pode registrar um fallback por haversine. Esse fallback é uma aproximação de triagem e não uma rota navegável confirmada.

\subsection{Agregação de emissões, custos e operações portuárias}
\label{subsec:agregacao}

Para cada alternativa $a$, as emissões operacionais e o proxy de custo são somas dos componentes incluídos:

\begin{equation}
E_a=\sum_{\ell\in L_a}E_{\ell},
\qquad
C_a=\sum_{\ell\in L_a}C_{\ell},
\label{eq:totais}
\end{equation}

em que $L_a$ é o conjunto de pernas e operações representadas. O consumo marítimo calculado pela Equação~\ref{eq:combustivel-cenario} é convertido em custo energético e emissões operacionais mediante os preços e fatores registrados no cenário. O custo não inclui margens comerciais, contratos, seguro, estoques, frequência ou tarifas finais de mercado.

Operações portuárias são representadas separadamente quando existe uma fonte utilizável, identificada como observação, média portuária, padrão de literatura ou indisponibilidade. Dado ausente não é convertido silenciosamente em zero. Quando a intensidade por trabalho de transporte do MRV é aplicada à navegação, o consumo anual reportado é tratado como abrangente para essa fronteira, e o hoteling do navio não é adicionado novamente, evitando dupla contagem. Componentes explícitos de terminal permanecem separados \cite{berth2009,shipops2022}.

\section{Implementação computacional}
\label{sec:implementacao}

O CabotageLens é implementado como aplicação Streamlit com lógica de domínio organizada em módulos reutilizáveis. A interface recebe origem, destino, carga e parâmetros de cenário. O backend resolve coordenadas, consulta rotas terrestres ou seus caches, seleciona portos, lê a matriz marítima enriquecida, calcula cada perna e apresenta totais, decomposição e proveniência.

O processamento marítimo é executado fora da interação do usuário. Os dados ANTAQ são normalizados em tabelas de viagens, paradas e chamadas; os workbooks anuais do EU MRV são reduzidos a um lookup por IMO e a estatísticas robustas por classe e tipo; por fim, a matriz marítima recebe os pares direcionais, alternativas de corredor, subtrechos e indicadores de cobertura. O artefato resultante preserva campos específicos para a intensidade representativa do par, para a intensidade observada do corredor selecionado e para as contagens de fonte.

No tempo de execução, a seleção da geometria e a avaliação econômica e ambiental utilizam esse artefato. A persistência em Supabase/Postgres armazena lugares, rotas e resultados reutilizáveis. A existência de cache reduz chamadas redundantes a provedores, mas não transforma uma rota modelada em trajetória GPS ou serviço contratado. O código-fonte e os artefatos rastreados estão disponíveis no repositório do projeto \cite{cabotagelensrepo}, e a aplicação é disponibilizada em interface web \cite{cabotagelensapp}.

\section{Evidência empírica e resultados}
\label{sec:resultados}

\subsection{Cobertura da base ANTAQ--EU MRV}

A base processada contém 1.324 viagens, 6.797 paradas, 7.103 ocorrências de chamadas portuárias e 389 IMOs únicos no recorte de cabotagem conteinerizada de 2025. O cruzamento com indicadores positivos do EU MRV associa 243 IMOs e 788 viagens de forma exata. Isso corresponde a 62,5\% dos IMOs, 59,5\% das viagens, 52,9\% da massa carregada e 50,6\% dos TEU observados. As viagens restantes recebem fallback identificado quando há classe ou tipo utilizável.

\begin{table}[htbp]
\centering
\small
\caption{Cobertura do cruzamento entre viagens ANTAQ e intensidade EU MRV.}
\label{tab:cobertura-mrv}
\begin{tabular}{lrr}
\toprule
Indicador & Valor & Cobertura \tabularnewline
\midrule
IMOs com correspondência exata & 243 de 389 & 62,5\% \tabularnewline
Viagens com correspondência exata & 788 de 1.324 & 59,5\% \tabularnewline
Carga em massa com correspondência exata & 15.959.761,561 de 30.191.845,948~t & 52,9\% \tabularnewline
Carga em TEU com correspondência exata & 1.454.351,75 de 2.872.715,00~TEU & 50,6\% \tabularnewline
\bottomrule
\end{tabular}
\end{table}

O enriquecimento da matriz marítima produz 177 pares direcionais e preserva 12.025 observações viagem--origem--destino. Entre 5.438 subtrechos canônicos calculados, 5.023 utilizam distância da matriz marítima e 415 registram fallback haversine. A intensidade é resolvida em 3.290 subtrechos por IMO e em 2.148 por fallback. A cobertura numérica de 100\% não deve ser confundida com cobertura observacional por navio: fontes exatas e imputadas permanecem separadas.

\subsection{Aplicação ao par Santos--Manaus}

O par Santos--Manaus contém 89 observações viagem--origem--destino com trabalho de transporte positivo, distribuídas em 22 sequências portuárias distintas. A população inclui a ligação direta e corredores com escalas intermediárias por diferentes portos; não há requisito de passagem por Suape ou por qualquer outro porto predeterminado. Das 89 observações, 40 utilizam intensidade exata por IMO e 49 utilizam fallback de tipo. Essas 49 observações concentram 59,54\% do trabalho de transporte do par.

A mediana ponderada definida pela Equação~\ref{eq:mediana-ponderada} é 9,322050~g/(t$\cdot$nm). O valor coincide com o fallback robusto de \textit{container ship}, pois as observações que usam essa fonte atravessam o percentil ponderado de 50\%. Para esse tipo, a amostra contém 243 valores positivos por IMO; a regra de 1\% remove dois valores de cada cauda e retém 239 observações. A média aritmética sem aparagem é 21,661852~g/(t$\cdot$nm), enquanto a mediana simples da amostra bruta é 4,620000~g/(t$\cdot$nm). Esses valores são diagnósticos da distribuição e não substituem o estimador registrado.

A geometria utilizada no cenário Santos--Manaus é o corredor direto observado, porque ele satisfaz o critério de seleção do caminho. Essa escolha não restringe a amostra usada na intensidade: as 89 viagens dos 22 corredores permanecem no cálculo do valor representativo. O resultado ilustra a separação metodológica entre o caminho empregado para somar distâncias e a população utilizada para estimar o desempenho energético do par.

\subsection{Benchmark externo}

O workbook associado aos estudos de Gustavo Costa fornece um conjunto documental de pares e parâmetros operacionais \cite{workbookdados}. Sua normalização identifica 21 pares comparáveis entre seis cidades para a unidade de 1~contêiner e 14~t. No próprio workbook, as emissões semanais agregadas são 7.614,97~tCO$_2$e no cenário rodoviário e 4.159,79~tCO$_2$e no cenário de cabotagem, diferença de aproximadamente 45,4\%. Esses números constituem referência externa de direção e ordem de grandeza; não calibram a intensidade marítima do CabotageLens nem eliminam diferenças de fatores, fronteiras e rotas.

\section{Discussão e limitações}
\label{sec:discussao}

A formulação marítima combina três níveis de evidência. A sequência portuária, os movimentos de carga e o IMO provêm da atividade registrada pela ANTAQ; a intensidade específica provém do EU MRV quando existe correspondência; e a falta de correspondência é preenchida por uma estatística robusta explicitamente identificada. Essa separação permite rastrear o que foi observado, o que foi derivado e o que foi imputado.

A reconstrução por subtrecho evita supor que o navio realizou somente a ligação analisada. Uma viagem Santos--Manaus com escalas intermediárias contribui com a carga efetivamente reconstruída e com a distância de cada trecho entre os dois portos. Ao mesmo tempo, o estimador do par aceita viagens de outros corredores quando compartilham a mesma origem e o mesmo destino na ordem correta. O peso por trabalho de transporte evita que uma viagem de baixa atividade tenha a mesma influência que uma viagem que transportou mais carga por maior distância.

As principais limitações são:

\begin{itemize}
  \item \textbf{Cobertura temporal e observabilidade:} o recorte de 2025 não representa automaticamente outros anos. Primeiras e últimas viagens da janela podem ser parciais, e chamadas físicas sem movimentação conteinerizada observável podem não aparecer na sequência normalizada.
  \item \textbf{Cobertura MRV parcial:} a resolução exata por IMO abrange parte da frota. Um fallback de classe ou tipo é uma imputação coletiva e não uma medição do navio ausente. No caso Santos--Manaus, a predominância do fallback no trabalho de transporte faz essa fonte controlar a mediana.
  \item \textbf{Tratamento de outliers:} a aparagem de 1\% protege o fallback contra caudas extremas, mas valores exatos por IMO são preservados. A mediana ponderada limita sua influência enquanto eles não concentram metade do trabalho de transporte; essa condição deve ser verificada em cada par.
  \item \textbf{Distâncias marítimas:} valores da matriz representam caminhos modelados. O fallback haversine é apenas uma aproximação de triagem e pode afetar a seleção do corredor mais curto.
  \item \textbf{Oferta de serviço:} uma sequência observada demonstra ocorrência na janela de dados, não frequência futura, disponibilidade de slot, regularidade comercial ou adequação do terminal à remessa analisada.
  \item \textbf{Fronteira ambiental:} os resultados representam emissões operacionais \ttw{} de \cooe{}. Não incluem produção e distribuição de combustíveis, construção de infraestrutura e ativos ou uma avaliação completa de ciclo de vida.
  \item \textbf{Fronteira econômica:} o valor monetário é um proxy de custo operacional modelado, não frete comercial, cotação de armador, contrato ou análise completa de competitividade.
\end{itemize}

Essas limitações impedem interpretar um resultado favorável como superioridade universal da cabotagem. A ferramenta oferece triagem comparativa e uma trilha de auditoria para o cenário modelado; decisões logísticas exigem evidência adicional de serviço, tempo, capacidade, terminal e preço comercial \cite{competitiveness2024,modalshiftreview2020}.

\section{Conclusão e trabalhos futuros}
\label{sec:conclusoes}

O CabotageLens estrutura uma comparação porta a porta entre transporte rodoviário direto e cadeia multimodal com cabotagem. Sua formulação marítima utiliza viagens observadas na ANTAQ, carga reconstruída em cada subtrecho e intensidades EU MRV associadas pelo IMO. Quando a associação exata não existe, uma estatística robusta de classe ou tipo é aplicada e identificada como fallback.

Para cada par direcional, todas as viagens nas quais a origem precede o destino participam do estimador, incluindo ligações diretas, múltiplas escalas e diferentes corredores. A intensidade representativa é a mediana das intensidades de viagem ponderada pelo trabalho de transporte. O corredor escolhido determina somente o caminho e a distância; dessa forma, a geometria do cenário não elimina evidências provenientes de outras sequências observadas para o mesmo par.

O par Santos--Manaus demonstra a aplicação do método: 89 observações em 22 corredores produzem intensidade de 9,322050~g/(t$\cdot$nm), com proveniência separada entre 40 observações por IMO e 49 por fallback. O resultado é auditável porque a matriz preserva amostra, pesos, fontes, alternativas de corredor e intensidade aplicada.

Trabalhos futuros devem ampliar a janela temporal da ANTAQ, aumentar a cobertura por IMO, incorporar oferta regular de serviços e refinar as distâncias marítimas. Também são extensões relevantes a representação portuária baseada em atividade observada, a análise de incerteza por fonte, a inclusão de preços comerciais verificáveis e a expansão da fronteira ambiental para uma avaliação \wtw{} ou de ciclo de vida.

\clearpage

\printbibliography[title={Referências}]

\end{document}
